## 4. Evaluation Metrics

A central contribution of our benchmark is identifying the limitations of standard personalization metrics when applied to multi-subject composition. Existing protocols heavily rely on global CLIP embeddings [3]. However, global embeddings are highly tolerant of local structural distortions and identity bleeding, rendering them inadequate for measuring severe identity entanglement. To provide a rigorous and comprehensive assessment that aligns with the P13N workshop's call for advanced evaluation metrics for personalization quality, we introduce a multi-tiered evaluation framework.

### 4.1. Text-Image Alignment (CLIP-T)
To measure how well the generated image aligns with the semantic intent of the text prompt, we compute the cosine similarity between the CLIP [3] text embedding of the prompt and the CLIP image embedding of the generated output. 

While CLIP-T is a standard metric for semantic fidelity, we observe a critical caveat in multi-subject scenarios: models may achieve high CLIP-T scores by generating a generic group of people that matches the macro-semantics of the prompt (e.g., "a group of people"), while completely failing to preserve the specific identities requested. Therefore, CLIP-T must be analyzed in conjunction with fine-grained identity metrics.

### 4.2. Identity Preservation: From CLIP to DINOv2
Traditionally, identity preservation is measured by computing the cosine similarity between the CLIP image embedding of the generated subject and the reference image (denoted as **CLIP-I**). However, our experiments reveal that CLIP-I scores remain artificially high even when subjects undergo severe identity bleeding or facial distortion.

To capture these local structural failures, we shift our primary identity evaluation to **DINOv2** [4], a self-supervised vision transformer known for its exceptional sensitivity to local features, part-level correspondence, and structural geometry. For each generated image $I_{gen}$ containing $N$ subjects, we compute the DINOv2 image embedding and calculate the average cosine similarity against the embeddings of all $N$ reference images $\{I_{ref}^1, I_{ref}^2, ..., I_{ref}^N\}$:

$$ \text{DINOv2 Score} = \frac{1}{N} \sum_{i=1}^{N} \cos(\text{DINOv2}(I_{gen}), \text{DINOv2}(I_{ref}^i)) $$

This metric effectively penalizes models that suffer from attention leakage, as DINOv2 strictly requires structural and identity consistency rather than mere stylistic or global semantic resemblance.

### 4.3. Subject Collapse Rate (SCR)
Average similarity scores can obscure catastrophic failures of individual subjects within a complex scene. For instance, in an 8-subject image, if 7 subjects are perfectly generated but 1 subject is completely missing or morphed into another identity, the mean DINOv2 score might still appear acceptable. 

<!-- 
=========================================
FIGURE 3 PLACEHOLDER (SCR METRIC EXPLANATION)
=========================================
\begin{figure}[t]
  \centering
  \includegraphics[width=\linewidth]{figures/scr_illustration.pdf}
  \caption{\textbf{Subject Collapse Rate (SCR).} Unlike average similarity scores which mask individual failures, SCR explicitly counts the proportion of subjects whose DINOv2 identity similarity falls below a strict threshold $\tau$. This provides a more realistic measure of multi-subject entanglement.}
  \label{fig:scr_illustration}
\end{figure}
-->

To explicitly quantify these localized failures, we propose the **Subject Collapse Rate (SCR)**. We define a subject as "collapsed" if its DINOv2 cosine similarity with the reference image falls below a predefined threshold $\tau$. The SCR for a given generated image is defined as the ratio of collapsed subjects to the total number of subjects:

$$ \text{SCR}_{@\tau} = \frac{\sum_{i=1}^{N} \mathbb{1}[\cos(\text{DINOv2}(I_{gen}), \text{DINOv2}(I_{ref}^i)) < \tau]}{N} $$

where $\mathbb{1}[\cdot]$ is the indicator function. Because DINOv2 similarities typically occupy a lower and more discriminative numerical range than CLIP, we employ strict thresholds $\tau \in \{0.4, 0.5, 0.6\}$. A lower SCR indicates better multi-subject preservation, while an SCR approaching 1.0 signifies a complete collapse of personalization.
